\begin{definition}[Dual Realization]\label{def:dualrealization}
    For a PIE system with corresponding PI operators $G\{\mcl{T}, \mcl{A}, \mcl{B}, \mcl{C}, \mcl{D}\}$ the state space representation is defined as
    \begin{equation}
        \bmat{\mcl{T}\dot{\mbf{x}}(t)\\ y(t)}
        =
        \bmat{\mcl{A} & \mcl{B}\\
              \mcl{C} & \mcl{D}}
        \bmat{\mbf{x}(t)\\ u(t)},
    \end{equation}
    define the dual realization, Theorem~\ref{thm:ogduality}, by
    \begin{equation}\label{eqn:dualrealization}
        G^\dual
        :=
        G\{\mcl{T}^*, \mcl{A}^*, \mcl{C}^*, \mcl{B}^*, \mcl{D}^*\}.
    \end{equation}
\end{definition}

\begin{definition}\label{def:primalscaledsys:iqc}
    For $w(t)\in\RL^{\mbf{n}_w}$, $\mbf{x}(t)\in\RL^{\mbf{n}_\mbf{x}}$, and 
    $z(t)\in\RL^{\mbf{n}_z}$ that satisfy $\mbf{x}(0) = 0$, the \textbf{Primal System} 
    $G\{\mcl{T}, \mcl{A}, \mcl{B}, \mcl{C}, \mcl{D}\}$ is defined as,
    \begin{equation}\label{eqn:primalsys:iqc}
        \bmat{\mcl{T}\dot{\mbf{x}}(t)\\ z(t)} = \bmat{\mcl{A} & \mcl{B} \\ 
        \mcl{C} & \mcl{D}}\bmat{\mbf{x}(t) \\ w(t)}
    \end{equation}
    where $\mcl{T}, \mcl{A}\in \Pi_4^{\mbf{n}_\mbf{x} \times \mbf{n}_\mbf{x}}$, 
    $\mcl{B}\in\Pi_4^{\mbf{n}_\mbf{x} \times \mbf{n}_{w}}$, 
    $\mcl{C}\in\Pi_4^{\mbf{n}_z \times \mbf{n}_\mbf{x}}$, and 
    $\mcl{D}\in\Pi_4^{\mbf{n}_z \times \mbf{n}_{w}}$ are PI operators.
    The \textbf{Scaled Primal System} is expressed in terms of $\tilde{w}(t)\in
    \RL^{\mbf{n}_w}$ and $\tilde{z}(t)\in\RL^{\mbf{n}_z}$, which satisfy,
    \begin{equation}\label{eqn:primalscaling:iqc}
        \Psi \bmat{z(t) \\ w(t)} = {\bmat{\Psi_{z} & \Psi_{zw} \\ 
        \Psi_{wz} & \Psi_{w}}}\bmat{z(t) \\ w(t)} = \bmat{\tilde{z}(t) \\ \tilde{w}(t)}.
    \end{equation}
    The operator $\Psi\in\mcl{L}(\RL)$ is called the \textbf{Primal Scaling} of the system.
    The state space representation of \eqref{eqn:primalscaling:iqc} is given by
    \begin{equation}\label{eqn:ssprimalscaling:iqc}
        \bmat{\mcl{T}_\Psi \mbf{x}_\Psi(t)\\ \tilde{z}(t) \\ \tilde{w}(t)}
        =
        \bmat{\mcl{A}_\Psi & \mcl{B}_\Psi^z & \mcl{B}_\Psi^w \\
            \mcl{C}_\Psi^z & \mcl{D}_\Psi^z & \mcl{D}_\Psi^{zw}\\
            \mcl{C}_\Psi^w & \mcl{D}_\Psi^{wz} & \mcl{D}_\Psi^w}\bmat{\mbf{x}_\Psi(t) \\ {z}(t) \\ {w}(t)}
    \end{equation}
    where the entries are PI operators of appropriate dimension.
\end{definition}



\begin{definition}\label{def:dualscalingoperation}
    Let $\Psi=\bmat{\Psi_z & \Psi_{zw}\\ \Psi_{wz} & \Psi_w}$ be an
    invertible scaling with state-space representation \eqref{eqn:ssprimalscaling:iqc}, and let
    $\Psi^\dual$ be defined by the dual realization. Let
    $\hat{\mcl{J}}:=\bmat{0&I\\ -I&0}$. Define the dual scaling operation
    \begin{equation}\label{eqn:dualoperation}
        \mbf{D}(\Psi)
        :=
        \left(\hat{\mcl{J}}^*\Psi^\dual\hat{\mcl{J}}\right)^{-1}
        =
        \bmat{\Psi_w^\dual & -\Psi_{zw}^\dual\\
              -\Psi_{wz}^\dual & \Psi_z^\dual}^{-1}.
    \end{equation}
    Since $\hat{\mcl{J}}^*=\hat{\mcl{J}}^{-1}$, $\hat{\mcl{J}}^* = -\hat{\mcl{J}}$, $\left(\Psi^{\dual}\right)^\dual = \Psi$, and $\left(\Psi^{-1}\right)^{-1}=\Psi$.
    Then, for any factorization $\Psi=UV$ in which $U$ and $V$ are invertible, the dual scaling operation is multiplicative,
    \begin{align}
        \mbf{D}(UV)
        &=
        \left(\hat{\mcl{J}}^*V^\dual U^\dual\hat{\mcl{J}}\right)^{-1} \nonumber\\
        &=
        \left(
        \hat{\mcl{J}}^*V^\dual\hat{\mcl{J}}
        \hat{\mcl{J}}^*U^\dual\hat{\mcl{J}}
        \right)^{-1} \nonumber\\
        &=
        \mbf{D}(U)\mbf{D}(V),
    \end{align}
    and the dual scaling operation is an involution,
    \begin{equation}\label{eqn:dualoperationinvolution}
        \mbf{D}(\mbf{D}(\Psi))=\Psi.
    \end{equation}
\end{definition}

\begin{definition}\label{def:dualscaledsys:iqc}
    For $\underline{w}(t)\in\RL^{\mbf{n}_z}$, $\underline{\mbf{x}}(t)\in\RL^{\mbf{n}_\mbf{x}}$, and 
    $\underline{z}(t)\in\RL^{\mbf{n}_w}$ that satisfy $\underline{\mbf{x}}(0) = 0$, the \textbf{Dual 
    System} $G^\dual$ is defined by applying the dual realization to the \textbf{Primal System}:
    \begin{equation}\label{eqn:dualsys:iqc}
        \bmat{\mcl{T}^*\dot{\underline{\mbf{x}}}(t)\\ \underline{z}(t)}
        =
        \bmat{\mcl{A}^* & \mcl{C}^*\\
              \mcl{B}^* & \mcl{D}^*}
        \bmat{\underline{\mbf{x}}(t) \\ \underline{w}(t)}.
    \end{equation}
    The \textbf{Scaled Dual System} is expressed in terms of $\tilde{\underline{w}}(t)\in
    \RL^{\mbf{n}_z}$ and $\tilde{\underline{z}}(t)\in\RL^{\mbf{n}_w}$, which satisfy,
    \begin{equation}\label{eqn:dualscaling:iqc}
        \mbf{D}(\Psi)
        \bmat{\underline{z}(t) \\ \underline{w}(t)} = 
        \bmat{\tilde{\underline{z}}(t) \\ \tilde{\underline{w}}(t)}.
    \end{equation}
    The operator $\mbf{D}(\Psi)\in\mcl{L}(\RL)$ is called the \textbf{Dual Scaling} of the system 
    and is the dual of the \textbf{Primal Scaling}.
    Thus the state space representation of \eqref{eqn:dualscaling:iqc} is given by
    \begin{equation}\label{eqn:ssdualscaling:iqc}
        \bmat{\mcl{T}_{\Psi}^*\dot{\underline{\mbf{x}}}_{\Psi}(t)\\ \underline{z}(t) \\ \underline{w}(t)}
        =
        \bmat{\mcl{A}_\Psi^* & -\mcl{C}_\Psi^{w*} & \mcl{C}_\Psi^{z*} \\
            -\mcl{B}_\Psi^{w*} & \mcl{D}_\Psi^{w*} & -\mcl{D}_\Psi^{zw*} \\
            \mcl{B}_\Psi^{z*} & -\mcl{D}_\Psi^{wz^*} & \mcl{D}_\Psi^{z*}}
            \bmat{\underline{\mbf{x}}_\Psi(t) \\ \underline{\tilde{z}}(t) \\ \underline{\tilde{w}}(t)}
    \end{equation}
\end{definition}

Adopted from \cite{venkataraman}
\begin{definition}(Potapov--Ginzburg Transform)\label{def:pg}
    Let ${\Psi} : \bmat{z \\ w} \to \bmat{\tilde{z} \\ \tilde{w}}$, where $\Psi = \bmat{\Psi_{z} & \Psi_{zw} \\ \Psi_{wz} & \Psi_{w}}$
    and $\Psi_{w}$ is invertible. Then the Potapov-Ginzburg transform of $\Psi$ is defined as
    \begin{equation}\label{eqn:pgtransform}
        \bmat{\tilde{z} \\ w}= \Psi^\dagger \bmat{\tilde{w} \\ z} \text{, where }\Psi^\dagger = \bmat{\Psi_{zw}\Psi_{w}^{-1} & \Psi_{z} - \Psi_{zw}\Psi_{w}^{-1}\Psi_{wz}\\
                            \Psi_{w}^{-1} & -\Psi_{w}^{-1}\Psi_{wz}}.
    \end{equation}
    Assumed to admit the state-space representation
    \begin{equation}\label{eqn:sspgtransform}
        \bmat{\mcl{T}_\Psi\dot{\mbf{x}}_\Psi(t)\\ \tilde{z}(t) \\ w(t)} = 
        \bmat{\hat{\mcl{A}}_\Psi & \hat{\mcl{B}}_\Psi^{\tilde{w}} & \hat{\mcl{B}}_\Psi^z\\
        \hat{\mcl{C}}_\Psi^z & \hat{\mcl{D}}_\Psi^{z\tilde{w}} & \hat{\mcl{D}}_\Psi^z \\
        \hat{\mcl{C}}_\Psi^w & \hat{\mcl{D}}_\Psi^{w\tilde{w}} & \hat{\mcl{D}}_\Psi^{wz}}\bmat{\mbf{x}_\Psi(t) \\ \tilde{w}(t) \\ z(t)},
    \end{equation}
    where
    \begin{align*}
        \hat{\mcl{A}}_\Psi &:= \mcl{A}_\Psi - \mcl{B}_\Psi^w(\mcl{D}_\Psi^w)^{-1}\mcl{C}_\Psi^w, &
        \hat{\mcl{B}}_\Psi^{\tilde{w}} &:= \mcl{B}_\Psi^w(\mcl{D}_\Psi^w)^{-1}, &
        \hat{\mcl{B}}_\Psi^{z} &:= \mcl{B}_\Psi^z - \mcl{B}_\Psi^w(\mcl{D}_\Psi^w)^{-1}\mcl{D}_\Psi^{wz}, &
        \hat{\mcl{C}}_\Psi^z &:= \mcl{C}_\Psi^z - \mcl{D}_\Psi^{zw}(\mcl{D}_\Psi^w)^{-1}\mcl{C}_\Psi^w,\\
        \hat{\mcl{D}}_\Psi^z &:= \mcl{D}_\Psi^z - \mcl{D}_\Psi^{zw}(\mcl{D}_\Psi^w)^{-1}\mcl{D}_\Psi^{wz}, &
        \hat{\mcl{D}}_\Psi^{z\tilde{w}} &:= \mcl{D}_\Psi^{zw}(\mcl{D}_\Psi^w)^{-1}, &
        \hat{\mcl{C}}_\Psi^w &:= -(\mcl{D}_\Psi^w)^{-1}\mcl{C}_\Psi^w, &
        \hat{\mcl{D}}_\Psi^{w\tilde{w}} &:= (\mcl{D}_\Psi^w)^{-1}, \\
        \hat{\mcl{D}}_\Psi^{wz} &:= -(\mcl{D}_\Psi^w)^{-1}\mcl{D}_\Psi^{wz}.
    \end{align*}
\end{definition}
\begin{lemma} (Potapov--Ginzburg Transform Equivalence)\label{lem:pgequivalence}
The \textbf{Scaled Primal System} (respectively, \textbf{Scaled Dual
System}) is equivalent, under an invertible congruence transformation,
to the \textbf{Primal System} (respectively, \textbf{Dual System})
augmented with the Potapov--Ginzburg transform of the \textbf{Primal
Scaling} (respectively, \textbf{Dual Scaling}).
\end{lemma}
\begin{proof}
    We first show the primal derivation explicitly. To show
    sufficiency, first augment the \textbf{Primal System} with the
    Potapov--Ginzburg transform of the \textbf{Primal Scaling}.
    Substitution of $z=\mcl{C}\mbf{x}+\mcl{D}w$ into
    \eqref{eqn:sspgtransform} gives the algebraic loop
    \begin{align*}
        w
            &=
            \Lambda_p\hat{\mcl{D}}_\Psi^{wz}\mcl{C}\mbf{x}
            +\Lambda_p\hat{\mcl{C}}_\Psi^w\mbf{x}_\Psi
            +\Lambda_p\hat{\mcl{D}}_\Psi^{w\tilde{w}}\tilde{w}.
    \end{align*}
    Assume that this loop is well-posed, and define
    \begin{equation*}
        \Lambda_p := (I-\hat{\mcl{D}}_\Psi^{wz}\mcl{D})^{-1},
        \qquad
        \Lambda_z := (I-\mcl{D}\hat{\mcl{D}}_\Psi^{wz})^{-1}.
    \end{equation*}
    Solving for $w$ and substituting into the plant state equation gives
    \begin{align*}
        \mcl{T}\dot{\mbf{x}}
            &= \left(\mcl{A}
                +\mcl{B}\Lambda_p\hat{\mcl{D}}_\Psi^{wz}\mcl{C}\right)\mbf{x}
               +\mcl{B}\Lambda_p\hat{\mcl{C}}_\Psi^w\mbf{x}_\Psi
               +\mcl{B}\Lambda_p\hat{\mcl{D}}_\Psi^{w\tilde{w}}\tilde{w}.
    \end{align*}
    Using $\mcl{D}\Lambda_p=\Lambda_z\mcl{D}$ and
    $I+\mcl{D}\Lambda_p\hat{\mcl{D}}_\Psi^{wz}=\Lambda_z$, as verified
    below in the push-through identities, \eqref{eqn:id1} and \eqref{eqn:id4}, we obtain
    \begin{equation*}
        z
        =
        \Lambda_z\mcl{C}\mbf{x}
        +\Lambda_z\mcl{D}\hat{\mcl{C}}_\Psi^w\mbf{x}_\Psi
        +\Lambda_z\mcl{D}\hat{\mcl{D}}_\Psi^{w\tilde{w}}\tilde{w}.
    \end{equation*}
    Substitution into the equations for $\dot{\mbf{x}}_\Psi$ and
    $\tilde{z}$ yields the augmented realization
    \begin{equation}\label{eqn:potapov}
    \bmat{\mcl{T}\dot{\mbf{x}}(t) \\ \mcl{T}_\Psi\dot{\mbf{x}}_\Psi(t)\\ \tilde{z}(t)} = 
    \bmat{\mcl{A} + \mcl{B}\Lambda_p\hat{\mcl{D}}_\Psi^{wz}\mcl{C} & \mcl{B}\Lambda_p\hat{\mcl{C}}_\Psi^w & \mcl{B}\Lambda_p\hat{\mcl{D}}_\Psi^{w\tilde{w}} \\
    \hat{\mcl{B}}_\Psi^z\Lambda_z\mcl{C} & \hat{\mcl{A}}_\Psi + \hat{\mcl{B}}_\Psi^z\Lambda_z\mcl{D}\hat{\mcl{C}}_\Psi^w & \hat{\mcl{B}}_\Psi^{\tilde{w}} + \hat{\mcl{B}}_\Psi^z\Lambda_z\mcl{D}\hat{\mcl{D}}_\Psi^{w\tilde{w}} \\
    \hat{\mcl{D}}_\Psi^z\Lambda_z\mcl{C} & \hat{\mcl{C}}_\Psi^z + \hat{\mcl{D}}_\Psi^z\Lambda_z\mcl{D}\hat{\mcl{C}}_\Psi^w & \hat{\mcl{D}}_\Psi^{z\tilde{w}} + \hat{\mcl{D}}_\Psi^z\Lambda_z\mcl{D}\hat{\mcl{D}}_\Psi^{w\tilde{w}}}\bmat{\mbf{x}(t) \\ \mbf{x}_\Psi(t) \\ \tilde{w}(t)}
    \end{equation}
Define also
\begin{equation*}
    \Lambda_c := (\mcl{D}_\Psi^{wz} \mcl{D} + \mcl{D}_\Psi^w)^{-1} = 
    (I+(\mcl{D}_\Psi^w)^{-1}\mcl{D}_\Psi^{wz} \mcl{D})^{-1}(\mcl{D}_\Psi^w)^{-1} = 
    (\mcl{D}_\Psi^w)^{-1} (I+\mcl{D}_\Psi^{wz}\mcl{D}(\mcl{D}_\Psi^w)^{-1})^{-1}.
\end{equation*}
The push-through identities
\begin{equation*}
    (I+\mcl{U}\mcl{V})^{-1}\mcl{U}
        = \mcl{U}(I+\mcl{V}\mcl{U})^{-1}, \qquad
    (I+\mcl{U}\mcl{V})^{-1}
        = I-\mcl{U}(I+\mcl{V}\mcl{U})^{-1}\mcl{V},
\end{equation*}
give the following relations
\begin{align}
    \Lambda_p(\mcl{D}_\Psi^w)^{-1}
        &= \left(I+(\mcl{D}_\Psi^w)^{-1}\mcl{D}_\Psi^{wz}\mcl{D}\right)^{-1}
           (\mcl{D}_\Psi^w)^{-1} = \Lambda_c
           \label{eqn:id1}\\[4pt]
    (\mcl{D}_\Psi^w)^{-1}\mcl{D}_\Psi^{wz}\Lambda_z
        &= (\mcl{D}_\Psi^w)^{-1}\mcl{D}_\Psi^{wz}
           \left(I+\mcl{D}(\mcl{D}_\Psi^w)^{-1}\mcl{D}_\Psi^{wz}\right)^{-1}
           \notag\\
        &= (\mcl{D}_\Psi^w)^{-1}
           \left(I+\mcl{D}_\Psi^{wz}\mcl{D}(\mcl{D}_\Psi^w)^{-1}\right)^{-1}
           \mcl{D}_\Psi^{wz}
         = \Lambda_c\mcl{D}_\Psi^{wz}
           \label{eqn:id2}\\[4pt]
    \Lambda_z\mcl{D}(\mcl{D}_\Psi^w)^{-1}
        &= \left(I+\mcl{D}(\mcl{D}_\Psi^w)^{-1}\mcl{D}_\Psi^{wz}\right)^{-1}
           \mcl{D}(\mcl{D}_\Psi^w)^{-1}
           \notag\\
        &= \mcl{D}(\mcl{D}_\Psi^w)^{-1}
           \left(I+\mcl{D}_\Psi^{wz}\mcl{D}(\mcl{D}_\Psi^w)^{-1}\right)^{-1}
         = \mcl{D}\Lambda_c
           \label{eqn:id3}\\[4pt]
    \Lambda_z
        &= \left(I+\mcl{D}(\mcl{D}_\Psi^w)^{-1}\mcl{D}_\Psi^{wz}\right)^{-1}
           \notag\\
        &= I-\mcl{D}(\mcl{D}_\Psi^w)^{-1}
           \left(I+\mcl{D}_\Psi^{wz}\mcl{D}(\mcl{D}_\Psi^w)^{-1}\right)^{-1}
           \mcl{D}_\Psi^{wz}
           \notag\\
        &= I-\mcl{D}\Lambda_c\mcl{D}_\Psi^{wz}
           \label{eqn:id4}\\[4pt]
    (\mcl{D}_\Psi^w)^{-1}
        -\Lambda_c\mcl{D}_\Psi^{wz}\mcl{D}(\mcl{D}_\Psi^w)^{-1}
        &= \left(I-\Lambda_c\mcl{D}_\Psi^{wz}\mcl{D}\right)
           (\mcl{D}_\Psi^w)^{-1}
           \notag\\
        &= \Lambda_c\left(\mcl{D}_\Psi^w
           +\mcl{D}_\Psi^{wz}\mcl{D}
           -\mcl{D}_\Psi^{wz}\mcl{D}\right)(\mcl{D}_\Psi^w)^{-1}
           \notag\\
        &= \Lambda_c.
           \label{eqn:id5}
\end{align}
We now express the blocks of \eqref{eqn:potapov} in terms of
$\Lambda_c$.  For the first row, direct substitution and
\eqref{eqn:id1} give
\begin{align*}
&\bmat{
    \mcl{A}+\mcl{B}\Lambda_p\hat{\mcl{D}}_\Psi^{wz}\mcl{C} &
    \mcl{B}\Lambda_p\hat{\mcl{C}}_\Psi^w &
    \mcl{B}\Lambda_p\hat{\mcl{D}}_\Psi^{w\tilde{w}}} =
\bmat{
    \mcl{A}-\mcl{B}\Lambda_c\mcl{D}_\Psi^{wz}\mcl{C} &
    -\mcl{B}\Lambda_c\mcl{C}_\Psi^w &
    \mcl{B}\Lambda_c}.
\end{align*}
For the second row,
\eqref{eqn:id2} and \eqref{eqn:id4} give
\begin{align*}
    \hat{\mcl{B}}_\Psi^z\Lambda_z
        &= \left(\mcl{B}_\Psi^z
            -\mcl{B}_\Psi^w(\mcl{D}_\Psi^w)^{-1}
             \mcl{D}_\Psi^{wz}\right)
            \Lambda_z\\
        &= \mcl{B}_\Psi^z\Lambda_z
            -\mcl{B}_\Psi^w(\mcl{D}_\Psi^w)^{-1}
             \mcl{D}_\Psi^{wz}\Lambda_z\\
        &= \mcl{B}_\Psi^z
            \left(I-\mcl{D}\Lambda_c\mcl{D}_\Psi^{wz}\right)
            -\mcl{B}_\Psi^w\Lambda_c\mcl{D}_\Psi^{wz}\\
        &= \mcl{B}_\Psi^z
            -\left(\mcl{B}_\Psi^z\mcl{D}+\mcl{B}_\Psi^w\right)
             \Lambda_c\mcl{D}_\Psi^{wz}.
\end{align*}
Its third entry follows from \eqref{eqn:id5}:
\begin{align*}
    &\hat{\mcl{B}}_\Psi^{\tilde{w}}
        +\hat{\mcl{B}}_\Psi^z\Lambda_z\mcl{D}
         \hat{\mcl{D}}_\Psi^{w\tilde{w}}\\
    &\quad =
        \mcl{B}_\Psi^w(\mcl{D}_\Psi^w)^{-1}
        +\left[
            \mcl{B}_\Psi^z
            -\left(\mcl{B}_\Psi^z\mcl{D}+\mcl{B}_\Psi^w\right)
             \Lambda_c\mcl{D}_\Psi^{wz}
         \right]\mcl{D}(\mcl{D}_\Psi^w)^{-1}\\
    &\quad =
        \left(\mcl{B}_\Psi^z\mcl{D}+\mcl{B}_\Psi^w\right)
        \left[
            (\mcl{D}_\Psi^w)^{-1}
            -\Lambda_c\mcl{D}_\Psi^{wz}\mcl{D}
             (\mcl{D}_\Psi^w)^{-1}
        \right]\\
    &\quad =
        \left(\mcl{B}_\Psi^z\mcl{D}+\mcl{B}_\Psi^w\right)\Lambda_c.
\end{align*}
The third row is obtained analogously:
\begin{equation*}
    \hat{\mcl{D}}_\Psi^z\Lambda_z
        = \mcl{D}_\Psi^z
           -\left(\mcl{D}_\Psi^z\mcl{D}+\mcl{D}_\Psi^{zw}\right)
            \Lambda_c\mcl{D}_\Psi^{wz} \qquad
    \hat{\mcl{D}}_\Psi^{z\tilde{w}}
        +\hat{\mcl{D}}_\Psi^z\Lambda_z\mcl{D}
         \hat{\mcl{D}}_\Psi^{w\tilde{w}}\\
    =
        \left(\mcl{D}_\Psi^z\mcl{D}+\mcl{D}_\Psi^{zw}\right)\Lambda_c.
\end{equation*}
Collecting these blocks and retaining $\tilde{w}$ as an output yields
\begin{equation}\label{eqn:congruence}
    \bmat{\mcl{T} \dot{\mbf{x}}(t) \\ \mcl{T}_\Psi\dot{\mbf{x}}_\Psi(t)\\ \tilde{z}(t) \\ \tilde{w}(t)} = 
    \bmat{\mcl{A} - \mcl{B}\Lambda_c\mcl{D}_\Psi^{wz}\mcl{C} & -\mcl{B}\Lambda_c\mcl{C}_\Psi^w & \mcl{B}\Lambda_c \\[4pt]
    \mcl{B}_\Psi^z \mcl{C} - (\mcl{B}_\Psi^z \mcl{D} + \mcl{B}_\Psi^w)\Lambda_c\mcl{D}_\Psi^{wz}\mcl{C} 
    & \mcl{A}_\Psi - (\mcl{B}_\Psi^z \mcl{D} + \mcl{B}_\Psi^w)\Lambda_c\mcl{C}_\Psi^w 
    & (\mcl{B}_\Psi^z \mcl{D} + \mcl{B}_\Psi^w)\Lambda_c \\[4pt]
    \mcl{D}_\Psi^z \mcl{C} - (\mcl{D}_\Psi^z \mcl{D} + \mcl{D}_\Psi^{zw})\Lambda_c\mcl{D}_\Psi^{wz}\mcl{C} 
    & \mcl{C}_\Psi^z - (\mcl{D}_\Psi^z \mcl{D} + \mcl{D}_\Psi^{zw})\Lambda_c\mcl{C}_\Psi^w 
    & (\mcl{D}_\Psi^z \mcl{D} + \mcl{D}_\Psi^{zw})\Lambda_c\\
    0&0&I}
    \bmat{\mbf{x}(t) \\ \mbf{x}_\Psi(t) \\ \tilde{w}(t)}
\end{equation}
The corresponding invertible congruence transformation is
\begin{equation}\label{eqn:congruencetf}
    \bmat{\mbf{x}(t) \\ \mbf{x}_\Psi(t) \\ \tilde{w}(t)} = 
    \bmat{I & 0 & 0\\ 
    0 & I & 0\\
    \mcl{D}_\Psi^{wz}\mcl{C} & \mcl{C}_\Psi^w & \Lambda_c^{-1}
    }\bmat{\mbf{x}(t) \\ \mbf{x}_\Psi(t) \\ {w}(t)}
    \qquad
    \bmat{\mbf{x}(t) \\ \mbf{x}_\Psi(t) \\ w(t)} = 
    \bmat{I & 0 & 0\\ 
    0 & I & 0\\
    -\Lambda_c\mcl{D}_\Psi^{wz} \mcl{C} & -\Lambda_c\mcl{C}_\Psi^w & \Lambda_c
    }\bmat{\mbf{x}(t) \\ \mbf{x}_\Psi(t) \\ \tilde{w}(t)}
\end{equation}
Applying \eqref{eqn:congruencetf} to \eqref{eqn:congruence} gives
\begin{equation}\label{eqn:congruence:transformed}
    \bmat{\mcl{T}\dot{\mbf{x}}(t)\\
          \mcl{T}_\Psi\dot{\mbf{x}}_\Psi(t)\\
          \tilde{z}(t)\\
          \tilde{w}(t)}
    =
    \bmat{
        \mcl{A} & 0 & \mcl{B}\\
        \mcl{B}_\Psi^z\mcl{C} & \mcl{A}_\Psi &
            \mcl{B}_\Psi^z\mcl{D}+\mcl{B}_\Psi^w\\
        \mcl{D}_\Psi^z\mcl{C} & \mcl{C}_\Psi^z &
            \mcl{D}_\Psi^z\mcl{D}+\mcl{D}_\Psi^{zw}\\
        \mcl{D}_\Psi^{wz}\mcl{C} & \mcl{C}_\Psi^w &
            \mcl{D}_\Psi^{wz}\mcl{D}+\mcl{D}_\Psi^w}
    \bmat{\mbf{x}(t)\\ \mbf{x}_\Psi(t)\\ w(t)}.
\end{equation}
The realization in \eqref{eqn:congruence:transformed} is the
direct augmentation of the \textbf{Primal System} with the \textbf{Primal Scaling}, and
hence is the \textbf{Scaled Primal System}. For necessity, applying \eqref{eqn:congruencetf} to
\eqref{eqn:congruence:transformed} recovers \eqref{eqn:congruence},
and reversing \eqref{eqn:id1}--\eqref{eqn:id5} recovers
\eqref{eqn:potapov}. The dual case follows by applying the same
argument to the \textbf{Dual System} and the \textbf{Dual Scaling}.
\end{proof}

\begin{theorem}[Scaled Dual $L_2$-gain]\label{thm:scaledduall2gain}
    Assume well-posedness. The following statements are equivalent.
    \begin{enumerate}
    \item For all $\tilde{w}\in\RL^{\mbf{n}_w}$ and $\mbf{x}(0)=0$, the \textbf{Scaled Primal System}
    satisfies $\|P_\tau\tilde{z}\| \leq \|P_\tau\tilde{w}\|$ for all $\tau\geq0$.
    \item For all $\tilde{\underline{w}}\in\RL^{\mbf{n}_z}$ and $\underline{\mbf{x}}(0)=0$, the \textbf{Scaled Dual System}
    satisfies $\|P_\tau\tilde{\underline{z}}\| \leq \|P_\tau\tilde{\underline{w}}\|$ for all $\tau\geq0$.
    \end{enumerate}
\end{theorem}
\begin{proof}
    The Potapov-Ginzburg transformation of \eqref{eqn:ssdualscaling:iqc} gives
    \begin{align*}
        \underline{\tilde{z}}
        &=
        (\mcl{D}_\Psi^{w*})^{-1}
            \mcl{B}_\Psi^{w*}\underline{\mbf{x}}_\Psi
        +(\mcl{D}_\Psi^{w*})^{-1}
            \mcl{D}_\Psi^{zw*}\underline{\tilde{w}}
        +(\mcl{D}_\Psi^{w*})^{-1}\underline{z}\\
        &=
        \bigl(\hat{\mcl{B}}_\Psi^{\tilde{w}}\bigr)^*
            \underline{\mbf{x}}_\Psi
        +\bigl(\hat{\mcl{D}}_\Psi^{z\tilde{w}}\bigr)^*
            \underline{\tilde{w}}
        +\bigl(\hat{\mcl{D}}_\Psi^{w\tilde{w}}\bigr)^*
            \underline{z},
    \end{align*}
    Substituting this expression into the
    first block row of \eqref{eqn:ssdualscaling:iqc} yields
    \begin{align*}
        \mcl{T}_\Psi^*\dot{\underline{\mbf{x}}}_\Psi
        &=
        \left(
            \mcl{A}_\Psi^*
            -
            \mcl{C}_\Psi^{w*}(\mcl{D}_\Psi^{w*})^{-1}
            \mcl{B}_\Psi^{w*}
        \right)\underline{\mbf{x}}_\Psi\\
        &\quad+
        \left(
            \mcl{C}_\Psi^{z*}
            -
            \mcl{C}_\Psi^{w*}(\mcl{D}_\Psi^{w*})^{-1}
            \mcl{D}_\Psi^{zw*}
        \right)\underline{\tilde{w}}
        -
        \mcl{C}_\Psi^{w*}(\mcl{D}_\Psi^{w*})^{-1}
            \underline{z}\\
        &=
        \bigl(\hat{\mcl{A}}_\Psi\bigr)^*
            \underline{\mbf{x}}_\Psi
        +
        \bigl(\hat{\mcl{C}}_\Psi^{z}\bigr)^*
            \underline{\tilde{w}}
        +
        \bigl(\hat{\mcl{C}}_\Psi^{w}\bigr)^*
            \underline{z}.
    \end{align*}
    Likewise, substituting into the third block row gives
    \begin{align*}
        \underline{w}
        &=
        \left(
            \mcl{B}_\Psi^{z*}
            -
            \mcl{D}_\Psi^{wz*}(\mcl{D}_\Psi^{w*})^{-1}
            \mcl{B}_\Psi^{w*}
        \right)\underline{\mbf{x}}_\Psi\\
        &\quad+
        \left(
            \mcl{D}_\Psi^{z*}
            -
            \mcl{D}_\Psi^{wz*}(\mcl{D}_\Psi^{w*})^{-1}
            \mcl{D}_\Psi^{zw*}
        \right)\underline{\tilde{w}}
        -
        \mcl{D}_\Psi^{wz*}(\mcl{D}_\Psi^{w*})^{-1}
            \underline{z}\\
        &=
        \bigl(\hat{\mcl{B}}_\Psi^{z}\bigr)^*
            \underline{\mbf{x}}_\Psi
        +
        \bigl(\hat{\mcl{D}}_\Psi^{z}\bigr)^*
            \underline{\tilde{w}}
        +
        \bigl(\hat{\mcl{D}}_\Psi^{wz}\bigr)^*
            \underline{z}.
    \end{align*}
    Collecting these three state-space relations gives
    \begin{equation}\label{eqn:dualpg}
    \bmat{\mcl{T}_\Psi^*\dot{\underline{\mbf{x}}}_\Psi(t) \\ \underline{\tilde{z}}(t) \\ \underline{w}(t)}
    =
    \bmat{
    \bigl(\hat{\mcl{A}}_\Psi\bigr)^*
    &
    \bigl(\hat{\mcl{C}}_\Psi^{z}\bigr)^*
    &
    \bigl(\hat{\mcl{C}}_\Psi^{w}\bigr)^*
    \\[2pt]
    
    \bigl(\hat{\mcl{B}}_\Psi^{\tilde{w}}\bigr)^*
    &
    \bigl(\hat{\mcl{D}}_\Psi^{z\tilde{w}}\bigr)^*
    &
    \bigl(\hat{\mcl{D}}_\Psi^{w\tilde{w}}\bigr)^*
    \\[2pt]
    
    \bigl(\hat{\mcl{B}}_\Psi^{z}\bigr)^*
    &
    \bigl(\hat{\mcl{D}}_\Psi^{z}\bigr)^*
    &
    \bigl(\hat{\mcl{D}}_\Psi^{wz}\bigr)^*
    }
    \bmat{\underline{\mbf{x}}_\Psi(t) \\ \underline{\tilde{w}}(t) \\ \underline{z}(t)}.
    \end{equation}
    We can now augment the dual system with \eqref{eqn:dualpg}.
    Substitution of
    $\underline{z}=\mcl{B}^*\underline{\mbf{x}}+\mcl{D}^*\underline{w}$
    into the last block row of \eqref{eqn:dualpg} gives the algebraic
    loop
    \begin{align*}
        \underline{w}
        &=
        \bigl(\hat{\mcl{D}}_\Psi^{wz}\bigr)^*
            \mcl{B}^*\underline{\mbf{x}}
        +\bigl(\hat{\mcl{B}}_\Psi^{z}\bigr)^*
            \underline{\mbf{x}}_\Psi
        +\bigl(\hat{\mcl{D}}_\Psi^{z}\bigr)^*
            \underline{\tilde{w}}
        +\bigl(\hat{\mcl{D}}_\Psi^{wz}\bigr)^*
            \mcl{D}^*\underline{w}.
    \end{align*}
    Assume this loop is well-posed. Since
    $\Lambda_z^*=(I-(\hat{\mcl{D}}_\Psi^{wz})^*\mcl{D}^*)^{-1}$,
    solving for $\underline{w}$ gives
    \begin{align*}
        \underline{w}
        &=
        \Lambda_z^*\bigl(\hat{\mcl{D}}_\Psi^{wz}\bigr)^*
            \mcl{B}^*\underline{\mbf{x}}
        +\Lambda_z^*\bigl(\hat{\mcl{B}}_\Psi^{z}\bigr)^*
            \underline{\mbf{x}}_\Psi
        +\Lambda_z^*\bigl(\hat{\mcl{D}}_\Psi^{z}\bigr)^*
            \underline{\tilde{w}}.
    \end{align*}
    Using
    $\mcl{D}^*\Lambda_z^*=\Lambda_p^*\mcl{D}^*$ and
    $I+\mcl{D}^*\Lambda_z^*(\hat{\mcl{D}}_\Psi^{wz})^*
    =\Lambda_p^*$, we also have
    \begin{align*}
        \underline{z}
        &=
        \Lambda_p^*\mcl{B}^*\underline{\mbf{x}}
        +\Lambda_p^*\mcl{D}^*
            \bigl(\hat{\mcl{B}}_\Psi^{z}\bigr)^*
            \underline{\mbf{x}}_\Psi
        +\Lambda_p^*\mcl{D}^*
            \bigl(\hat{\mcl{D}}_\Psi^{z}\bigr)^*
            \underline{\tilde{w}}.
    \end{align*}

    To prove sufficiency, suppose the first statement holds. By Lemma~\ref{lem:pgequivalence},
    the realization in \eqref{eqn:potapov} satisfies
    $\|P_\tau\tilde{z}\|\leq\|P_\tau\tilde{w}\|$ for all $\tau\geq0$.
    Applying the argument of Theorem~\ref{thm:ogduality} with
    $\gamma=1$ to \eqref{eqn:potapov} gives its dual realization 
    ($\mcl{D}\Lambda_p=\Lambda_z\mcl{D}$, $\Lambda_p\hat{\mcl{D}}_\Psi^{wz} = \hat{\mcl{D}}_\Psi^{wz}\Lambda_z$)
    \begin{equation}\label{eqn:dualpotapov}
    \bmat{\mcl{T}^*\dot{\underline{\mbf{x}}}(t) \\
          \mcl{T}_\Psi^*\dot{\underline{\mbf{x}}}_\Psi(t)\\
          \underline{\tilde{z}}(t)}
    =
    \bmat{
        \mcl{A}^*
            +\mcl{C}^*\Lambda_z^*
             \bigl(\hat{\mcl{D}}_\Psi^{wz}\bigr)^*\mcl{B}^*
        &
        \mcl{C}^*\Lambda_z^*
             \bigl(\hat{\mcl{B}}_\Psi^{z}\bigr)^*
        &
        \mcl{C}^*\Lambda_z^*
             \bigl(\hat{\mcl{D}}_\Psi^{z}\bigr)^*
        \\[4pt]
        \bigl(\hat{\mcl{C}}_\Psi^w\bigr)^*
            \Lambda_p^*\mcl{B}^*
        &
        \bigl(\hat{\mcl{A}}_\Psi\bigr)^*
            +\bigl(\hat{\mcl{C}}_\Psi^w\bigr)^*
             \Lambda_p^*\mcl{D}^*
             \bigl(\hat{\mcl{B}}_\Psi^{z}\bigr)^*
        &
        \bigl(\hat{\mcl{C}}_\Psi^z\bigr)^*
            +\bigl(\hat{\mcl{C}}_\Psi^w\bigr)^*
             \Lambda_p^*\mcl{D}^*
             \bigl(\hat{\mcl{D}}_\Psi^{z}\bigr)^*
        \\[4pt]
        \bigl(\hat{\mcl{D}}_\Psi^{w\tilde{w}}\bigr)^*
            \Lambda_p^*\mcl{B}^*
        &
        \bigl(\hat{\mcl{B}}_\Psi^{\tilde{w}}\bigr)^*
            +\bigl(\hat{\mcl{D}}_\Psi^{w\tilde{w}}\bigr)^*
             \Lambda_p^*\mcl{D}^*
             \bigl(\hat{\mcl{B}}_\Psi^{z}\bigr)^*
        &
        \bigl(\hat{\mcl{D}}_\Psi^{z\tilde{w}}\bigr)^*
            +\bigl(\hat{\mcl{D}}_\Psi^{w\tilde{w}}\bigr)^*
             \Lambda_p^*\mcl{D}^*
             \bigl(\hat{\mcl{D}}_\Psi^{z}\bigr)^*}
    \bmat{\underline{\mbf{x}}(t) \\
          \underline{\mbf{x}}_\Psi(t) \\
          \underline{\tilde{w}}(t)}.
\end{equation}
    satisfies
    $\|P_\tau\tilde{\underline{z}}\|\leq
    \|P_\tau\tilde{\underline{w}}\|$ for all $\tau\geq0$.
    By Lemma~\ref{lem:pgequivalence}. The realization in \eqref{eqn:dualpotapov} is the
    \textbf{Dual System} augmented with the Potapov--Ginzburg transform
    \eqref{eqn:dualpg} of the \textbf{Dual Scaling}. Hence by Lemma~\ref{lem:pgequivalence} this is equivalent to the
    \textbf{Scaled Dual System} satisfying the same inequality.
    Since $(G^\dual)^\dual=G$ and  $\mbf{D}(\mbf{D}(\Psi))=\Psi$ we have that sufficiency implies necessity.
\end{proof}


\begin{example}[Dual multiplier for the zero-padded LTI-parametric IQC]
    Consider an LTI-parametric IQC with causal bounded linear tall factor
    $\psi=\bmat{I\\\psi_t}$ and pointwise multiplier with $\mcl{R}=-\mcl{Q}$,
    \begin{equation}
        \mcl{P}:=
        \bmat{\mcl{Q} & \mcl{S}\\ \mcl{S}^* & -\mcl{Q}},
        \qquad
        \mcl{Q}\succ0,\quad
        \mcl{S}=-\mcl{S}^* .
    \end{equation}
    Let $\mcl{M}\in\mcl{L}(\mbf L_{[a,b]})$ be the multiplication operator
    induced by $\mcl{P}$, so, for $v\in\mbf L_{[a,b]}$,
    \begin{equation}
        (\mcl{M}v)(t):=\mcl{P}v(t).
    \end{equation}
    Define the original tall filtered map
    \begin{equation}
        \Psi_{\mathrm{tall}}
        :=
        \bmat{\psi & 0\\0 & \psi}.
    \end{equation}
    In this example we use the system-side convention this is the
    opposite sign from the uncertainty-side IQC convention in
    Definition~\ref{def:multipliers}. Assume the following inequality holds
    for all $T\geq0$ and all $z(t),w(t)\in\RL$ where ${z}(t)$ satisfies Definition~\ref{eqn:primalsys:iqc},,
    \begin{equation}
        \ip{
            p_T\Psi_{\mathrm{tall}}\bmat{z\\w}
        }{
            p_T\mcl{M}\Psi_{\mathrm{tall}}\bmat{z\\w}
        }_{\mbf L_{[a,b]}}
        \leq 0.
    \end{equation}
    Since Definition~\ref{def:dualscalingoperation} applies to square
    invertible scalings, introduce
    \begin{equation}
        \mcl{N}:=\bmat{I&0},
        \qquad
        \mcl{N}^*=\bmat{I\\0},
        \qquad
        \psi=\psi_\square\mcl{N}^*,
    \end{equation}
    \begin{equation}
        \psi_\square
        :=
        \bmat{I & 0 \\ \psi_t & I}
        \bmat{I & \psi_t^\dual\\ 0 & I}
        =
        \bmat{I & \psi_t^\dual \\
              \psi_t & I+\psi_t\psi_t^\dual}.
    \end{equation}
    By construction $\psi_t^\dual$ is causal bounded.
    Then $\psi_\square$ is square and invertible, with causal bounded inverse
    \begin{equation}
        \psi_\square^{-1}
        =
        \bmat{I & -\psi_t^\dual \\ 0 & I}
        \bmat{I & 0 \\ -\psi_t & I}
        =
        \bmat{
            I+\psi_t^\dual\psi_t & -\psi_t^\dual\\
            -\psi_t & I
        }.
    \end{equation}
    Square up the plant by setting $z_\square=\mcl{N}^*z$ and
    $w=\mcl{N}w_\square$:
    \begin{equation}
        \bmat{\mcl{T}\dot{\mbf{x}}(t)\\ z_\square(t)}
        =
        \bmat{
        \mcl{A} & \mcl{B}\mcl{N}\\
        \mcl{N}^*\mcl{C} & \mcl{N}^*\mcl{D}\mcl{N}
        }
        \bmat{\mbf{x}(t)\\ w_\square(t)}.
    \end{equation}
    The original plant is recovered on the embedded square subspace by taking
    $w_\square=\mcl{N}^*w=\bmat{w\\0}$. Hence, for
    $z_\square=\mcl{N}^*z=\bmat{z\\0}$ and
    $w_\square=\mcl{N}^*w=\bmat{w\\0}$, the squared-up primal J-contraction
    inequality is
    \begin{equation}
        \Psi_\square
        :=
        \bmat{\psi_\square&0\\0&\psi_\square},
        \qquad
        \ip{
            p_T\Psi_\square\bmat{z_\square\\w_\square}
        }{
            p_T\mcl{M}\Psi_\square\bmat{z_\square\\w_\square}
        }_{\mbf L_{[a,b]}}
        \leq0.
    \end{equation}
    Factor the pointwise multiplier as $\mcl{P}=\mcl{V}^*J\mcl{V}$ with
    \begin{equation}
        \mcl{V}
        :=
        \bmat{\mcl{V}_{11} & \mcl{V}_{12}\\0 & \mcl{V}_{22}},
        \qquad
        J:=
        \bmat{I & 0\\0 & -I},
    \end{equation}
    where $\mcl{V}_{11}$ and $\mcl{V}_{22}$ are invertible. Define the
    $J$-factorization
    \begin{equation}
        \Phi
        :=
        \mcl{V}\Psi_\square
        =
        \bmat{
            \mcl{V}_{11}\psi_\square
            &
            \mcl{V}_{12}\psi_\square\\
            0
            &
            \mcl{V}_{22}\psi_\square
        }.
    \end{equation}
    Using Definition~\ref{def:dualscalingoperation}
    gives the dual scaling
    \begin{equation}
        \mbf{D}(\Phi)
        =
        \mbf{D}(\mcl{V})\mbf{D}\!\left(
        \bmat{\psi_\square & 0\\0 & \psi_\square}
        \right)
        =
        \underbrace{
        \bmat{
            \mcl{V}_{22}^{-*}
            &
            \mcl{V}_{22}^{-*}\mcl{V}_{12}^*\mcl{V}_{11}^{-*}\\
            0
            &
            \mcl{V}_{11}^{-*}
        }}_{\mcl{V}_d}
        \underbrace{\bmat{
            \psi_\square^{-\dual} & 0\\
            0 & \psi_\square^{-\dual}
        }}_{\Psi_{d,\square}} .
    \end{equation}
    Hence the pointwise multiplication operator $\mcl{M}_d\in\mcl{L}(\mbf L_{[a,b]})$
    by
    \begin{equation}
        \mcl{P}_d:=\mcl{V}_d^*J\mcl{V}_d,
        \qquad
        (\mcl{M}_d v)(t):=\mcl{P}_d v(t).
    \end{equation}
    Thus the squared-up dual contraction inequality is
    \begin{equation}
        \ip{
            p_T\Psi_{d,\square}
            \bmat{\underline{z}_{\square}\\\underline{w}_{\square}}
        }{
            p_T\mcl{M}_d\Psi_{d,\square}
            \bmat{\underline{z}_{\square}\\\underline{w}_{\square}}
        }_{\mbf L_{[a,b]}}
        \leq0.
    \end{equation}
    Equivalently,
    \begin{equation}
        \mcl{P}_d
        =
        \bmat{\mcl{Q}_d & \mcl{S}_d\\
              \mcl{S}_d^* & \mcl{R}_d}
        =
        \mcl{K}^{-1},
        \qquad
        \mcl{K}
        :=
        \bmat{\mcl{Q} & \mcl{S}^*\\
              \mcl{S} & -\mcl{Q}} .
    \end{equation}
    With
    \begin{equation}
        \Theta
        :=
        \mcl{Q}+\mcl{S}^*\mcl{Q}^{-1}\mcl{S}
        =
        \mcl{Q}-\mcl{S}\mcl{Q}^{-1}\mcl{S},
    \end{equation}
    the block inverse is
    \begin{equation}
        \mcl{P}_d
        =
        \bmat{
            \Theta^{-1}
            &
            \Theta^{-1}\mcl{S}^*\mcl{Q}^{-1}\\
            \mcl{Q}^{-1}\mcl{S}\Theta^{-1}
            &
            -\Theta^{-1}
        },
    \end{equation}
    so
    \begin{equation}
        \mcl{Q}_d=\Theta^{-1}\succ0,
        \qquad
        \mcl{S}_d
        =
        \Theta^{-1}\mcl{S}^*\mcl{Q}^{-1}
        =
        -\Theta^{-1}\mcl{S}\mcl{Q}^{-1},
        \qquad
        \mcl{R}_d=-\Theta^{-1}=-\mcl{Q}_d\prec0.
    \end{equation}
    Moreover, since
    \begin{equation}
        \mcl{S}\mcl{Q}^{-1}\Theta
        =
        \Theta\mcl{Q}^{-1}\mcl{S}
        =
        \mcl{S}
        -
        \mcl{S}\mcl{Q}^{-1}\mcl{S}\mcl{Q}^{-1}\mcl{S},
    \end{equation}
    we have $\Theta^{-1}\mcl{S}\mcl{Q}^{-1}
    =\mcl{Q}^{-1}\mcl{S}\Theta^{-1}$, and therefore
    \begin{equation}
        \mcl{S}_d^*
        =
        \mcl{Q}^{-1}\mcl{S}\Theta^{-1}
        =
        \Theta^{-1}\mcl{S}\mcl{Q}^{-1}
        =
        -\mcl{S}_d,
        \qquad\text{equivalently}\qquad
        \mcl{S}_d=-\mcl{S}_d^* .
    \end{equation}
    To express the dual condition on the original signal dimension, restrict
    the squared-up dual variables to the subspace
    $\underline{z}_\square=\mcl{N}^*\underline{z}$ and
    $\underline{w}_\square=\mcl{N}^*\underline{w}$, so that
    $\underline{z}=\mcl{N}\underline{z}_\square$ and
    $\underline{w}=\mcl{N}\underline{w}_\square$. This gives
    \begin{equation}
        \psi_d
        :=
        \psi_\square^{-\dual}\mcl{N}^*
        =
        \bmat{I+\psi_t^\dual\psi_t\\-\psi_t}.
    \end{equation}
    Thus
    \begin{equation}
        \Psi_d = \bmat{\psi_d & 0\\0 & \psi_d}
        .
    \end{equation}
    The induced dual contraction inequality is therefore, for all
    $T\geq0$ and all $\underline{z}(t),\underline{w}(t)\in\RL$, where $\underline{z}(t)$ satisfies Definition~\ref{eqn:dualsys:iqc},
    \begin{equation}
        \ip{
            p_T
            \Psi_d
            \bmat{\underline{z}\\\underline{w}}
        }{
            p_T\mcl{M}_d
            \Psi_d
            \bmat{\underline{z}\\\underline{w}}
        }_{\mbf L_{[a,b]}}
        \leq0.
    \end{equation}
\end{example}
